% Options for packages loaded elsewhere
\PassOptionsToPackage{unicode}{hyperref}
\PassOptionsToPackage{hyphens}{url}
\documentclass[
]{article}
\usepackage{xcolor}
\usepackage[margin=1in]{geometry}
\usepackage{amsmath,amssymb}
\setcounter{secnumdepth}{-\maxdimen} % remove section numbering
\usepackage{iftex}
\ifPDFTeX
  \usepackage[T1]{fontenc}
  \usepackage[utf8]{inputenc}
  \usepackage{textcomp} % provide euro and other symbols
\else % if luatex or xetex
  \usepackage{unicode-math} % this also loads fontspec
  \defaultfontfeatures{Scale=MatchLowercase}
  \defaultfontfeatures[\rmfamily]{Ligatures=TeX,Scale=1}
\fi
\usepackage{lmodern}
\ifPDFTeX\else
  % xetex/luatex font selection
\fi
% Use upquote if available, for straight quotes in verbatim environments
\IfFileExists{upquote.sty}{\usepackage{upquote}}{}
\IfFileExists{microtype.sty}{% use microtype if available
  \usepackage[]{microtype}
  \UseMicrotypeSet[protrusion]{basicmath} % disable protrusion for tt fonts
}{}
\makeatletter
\@ifundefined{KOMAClassName}{% if non-KOMA class
  \IfFileExists{parskip.sty}{%
    \usepackage{parskip}
  }{% else
    \setlength{\parindent}{0pt}
    \setlength{\parskip}{6pt plus 2pt minus 1pt}}
}{% if KOMA class
  \KOMAoptions{parskip=half}}
\makeatother
\usepackage{graphicx}
\makeatletter
\newsavebox\pandoc@box
\newcommand*\pandocbounded[1]{% scales image to fit in text height/width
  \sbox\pandoc@box{#1}%
  \Gscale@div\@tempa{\textheight}{\dimexpr\ht\pandoc@box+\dp\pandoc@box\relax}%
  \Gscale@div\@tempb{\linewidth}{\wd\pandoc@box}%
  \ifdim\@tempb\p@<\@tempa\p@\let\@tempa\@tempb\fi% select the smaller of both
  \ifdim\@tempa\p@<\p@\scalebox{\@tempa}{\usebox\pandoc@box}%
  \else\usebox{\pandoc@box}%
  \fi%
}
% Set default figure placement to htbp
\def\fps@figure{htbp}
\makeatother
\setlength{\emergencystretch}{3em} % prevent overfull lines
\providecommand{\tightlist}{%
  \setlength{\itemsep}{0pt}\setlength{\parskip}{0pt}}
\usepackage{bookmark}
\IfFileExists{xurl.sty}{\usepackage{xurl}}{} % add URL line breaks if available
\urlstyle{same}
\hypersetup{
  pdftitle={part4-vera},
  hidelinks,
  pdfcreator={LaTeX via pandoc}}

\title{part4-vera}
\author{}
\date{\vspace{-2.5em}2026-04-27}

\begin{document}
\maketitle

\#4. Volatility Model with ARMA-GARCH \#\#4.1 Sample Period and Data
Description Following the project instructions, the analysis focuses on
the most recent market dynamics. We consider daily data from January
2020 to December 2024, corresponding to the last 5 years. This period
captures several relevant market regimes, including the COVID-19
financial shock, the post-pandemic recovery, the monetary tightening
cycles and the increased volatility in cryptocurrency and technology
sectors. Using five years of data provides us with a balance between
enough observations for an accurate estimation of conditional volatility
models and structural stability to avoid outdated market behaviour. Log
returns were computed as \[
r_t = \log(P_t) - \log(P_{t-1})
\] where Pt is the daily closing prices.

\#\#4.2 ARCH Effects \#\#\# Analysis of the ARMA Residual The ARMA(p,q)
model selected was estimated and the residuals were extracted, for each
time series and we tested for conditional heteroskedasticity using the
ARCH-LM test with 10 lags, to check whether volatility modelling is
required. Under the null hypothesis of no ARCH effects, the statistic
follows a chi-squared distribution with 10 degrees of freedom. The
results indicate that all assets exhibit a p-value \textless{} 0.05,
which indicates statistically significant ARCH effects, implying that
there is volatility clustering, conditional variance that varies over
time and inadequacy of constant-variance ARMA models.

\subsubsection{Diagnostics of the Squared
Residuals}\label{diagnostics-of-the-squared-residuals}

To analyze the squared residuals we used the Box-Pierce test on squared
residuals, revealing serial dependence and the autocorrelation in
squared residuals confirms that there is a persistence in volatility.
These results are consistent with the stylized facts of the financial
returns, such as, volatility clustering, heavy tails and the persistence
of shocks.

\subsubsection{Proposed ARMA-GARCH
Models}\label{proposed-arma-garch-models}

We propose modelling the log returns using ARMA(p,q) - GARCH(1,1). The
GARCH(1,1) model was selected to model the extended persistence of
volatility because it's a more parsimonious model to describe volatility
properties and because time series models should be as simple as
possible so it wouldn't be worth it to use a higher-order model. The
conditional variance is described by the GARCH as depending on both
lagged squared errors and conditional variances, with the equation: \[
\sigma_t^2 = \omega + \alpha \varepsilon_{t-1}^2 + \beta \sigma_{t-1}^2
\]

\subsection{4.3 Adequacy of the model and
diagnostics}\label{adequacy-of-the-model-and-diagnostics}

After estimating the models with normally distributed errors, we
performed three diagnostic tests:

\#\#\#Standardized Residual ACF The autocorrelation functions of the
standardized residuals show no significant remaining serial dependence,
which indicates that there is an adequate mean specification.
Significant spikes in the graph would indicate that the ARMA component
had not fully captured the mean dynamics.

GRAPH

\#\#\#Squared Standardized Residual ACF The absence of significant
autocorrelation shows that conditional heteroskedasticity has been
successfully captured.

GRAPH

\#\#\#Normal Q-Q Plots These plots reveal deviations in the tails for
many assets such as BTC-USD, SOL-USD and RKLB, suggesting excess
kurtosis compared to the Gaussian distribution.

GRAPH

\#\#\#Model Selection Overall, the ARMA-GARCH(1,1) model is the selected
under normality, since it captures the return dynamics, volatility
clustering and the persistence of conditional variance.

\subsection{4.3 ARMA-GARCH with Student-t
Errors}\label{arma-garch-with-student-t-errors}

We re-estimated the models assuming a Student-t distribution to account
for the heavy tails observed in the Q-Q plots. We obtained an
improvement in Q-Q plot alignment and reduced tail deviations, as for
the ACF diagnostics, they remain clean, confirming that serial
correlation and ARCH effects continue adequate. The Student-t
distribution accommodates extreme observations more effectively, which
are frequent in cryptocurrencies and the growth stocks included in the
sample.

\subsection{4.4 Forecasting Returns and
Volatility}\label{forecasting-returns-and-volatility}

\#\#\#Training and Test Sets Each time series was divided into a
training set with the first 1000 observations and a test set with the
remaining ones. The rolling one-step-ahead forecasts were generated
using GARCH(1,1) which is the parsimonious benchmark discussed
previously and GARCH(2,1) which is an extended specification that allows
the variance to depend on two lagged squared innovations, capturing more
complex volatility dynamics. The forecasts were produced for the
conditional mean (log returns) and the conditional volatility.

\#\#\#Volatility Forecasts With this tests, we discovered that both
GARCH(1,1) and GARCH(2,1) produce smooth volatility forecasts that are
able to capture the clustering of large shocks, during high-volatility
episodes, both models appropriately increase their volatility forecasts
with a short lag, GARCH(2,1) does not systematically outperform
GARCH(1,1) in terms of volatility tracking since for most assets both
forecasts are very similar and lastly, we discovered that cryptocurrency
assets exhibit the most volatile forecast path, with some large spikes
that neither model fully accounts for.

\#\#\#Log Return Forecasts The log return forecasts are, as expected,
close to the unconditional mean for most periods. The ARMA mean
component provides modest short-run predictability for assets with
non-zero parameters, but the returns are controlled by unpredictable
shocks.

Overall, GARCH(1,1) is preferred for parsimony because it provides
stable forecasts, however GARCH(2,1) might be used occasionally since it
captures longer volatility persistence but adds complexity.

\subsection{4.5 1st and 2nd Strep
Forecasts}\label{st-and-2nd-strep-forecasts}

We illustrated the multi-step forecast by selecting AMD (Advanced Micro
Devices), which has the ARMA(3,0)-GARCH(1,1)/GARCH(2/1) specification.
The two-step forecasts were obtained from both GARCH specifications and
computed as:

\#\#\#Step 1 Forecast ???? don't get it

\#\#\#Step 2 Forecast also don't get it

The table presented reports the 1st and 2nd step forecasts for the
predicted returns and volatility.

insert table

Results show convergence of the volatility forecasts into the long-run
variance.

\subsection{4.6 Economic Interpretation Tests (mean, leverage and risk
premia)}\label{economic-interpretation-tests-mean-leverage-and-risk-premia}

We did three disticnt hypothesis tests for each asset:

\subsubsection{1. Mean of Log Returns}\label{mean-of-log-returns}

Our null hypothesis was of mean zero and the majority of the assets show
insignificant mean parameters, which indicates that the average returns
are not statistically different from zero. This result is consistent
with the weak-form efficient hypothesis.

\#\#\#2. Leverage Effects To test the leverage effects, we used an
EGARCH(1,1) model where a significant negative asymmetry parameter
indicates that negative shocks increase volatility, more than positive
shocks. It does not converge for all assets, suggesting that asymmetric
volatility is asset-specific, for example, cryptocurrencies tend to show
less consistent leverage since they are a very speculative and sentiment
driven stocks.

\#\#\#3. Risk Premia For the risk premia, we used an GARCH-in-mean
specification, with the null hypothesis of no relationship between
expected return and risk. Only some assets showed significant
ARCH-in-mean coefficients, which implies that there is a weak evidence
that higher conditional volatility leads to higher expected returns.

Overall, our analysis confirms that the key stylized facts of financial
markets are the strong volatility clustering, the heavy-tailed return
distributions and the persistent conditional variance.

\end{document}
